\chapter{Background and Related Work}
\label{ch:background}

\section{Retrieval-augmented generation}
\label{sec:bg-rag}

A RAG system has three stages \cite{lewis2020rag}. An \emph{embedder} maps
every passage in the collection, and each incoming query, into a shared vector
space. Here the embedder is a sentence transformer \cite{reimers2019sbert}. A
\emph{retriever} returns the $k$ passages whose vectors are nearest the
query's. At the scale of hundreds of thousands of passages this is done
approximately, with a graph index such as HNSW \cite{malkov2020hnsw} in a
library such as FAISS \cite{johnson2019faiss}. A \emph{generator}, the LLM,
receives the query and the $k$ passages in a prompt and writes the answer.

Two properties of this pipeline matter for what follows. First, the retriever
ranks by similarity, not by truth or provenance, so a passage written to
resemble a query will be retrieved for it. Second, the generator sees only the
top $k$. If the top $k$ is filled with adversarial passages, the passage that
actually answers the query is \emph{evicted} and the model never sees it.

Production systems often add a second-stage \emph{reranker}: a cross-encoder
that scores each (query, passage) pair jointly and reorders a larger candidate
pool \cite{nogueira2019rerank}. Rerankers reliably improve retrieval quality
on clean queries. \Cref{sec:def-rerank} shows that this does not carry over to
adversarial ones.

\section{Alignment training and context-faithfulness}
\label{sec:bg-alignment}

A pretrained (\emph{base}) model continues text. An instruction-tuned
(\emph{aligned}) model is further trained on instruction--response pairs and,
usually, on human or AI preference data, so that it follows instructions and
behaves as an assistant \cite{wei2022flan,ouyang2022instructgpt}. All four
model families in this thesis release both versions from the same pretrained
checkpoint \cite{grattafiori2024llama3,qwen2025qwen25,jiang2023mistral}, which
is what makes a matched comparison possible.

Alignment changes refusal behaviour in two distinct ways. \emph{Safety}
training teaches a model to decline harmful requests. Its failure mode,
refusing benign requests that superficially resemble harmful ones, is measured
by over-refusal benchmarks \cite{rottger2024xstest,cui2024orbench}.
\emph{Context-faithfulness} is a different property. It is the tendency to
answer from the supplied context rather than from parametric memory, and to
abstain when the context does not support an answer
\cite{longpre2021entity,zhou2023contextfaithful,xie2024chameleon}. In RAG it is
usually treated as desirable, since it is the main defence against
hallucination. This thesis argues that it is also the attack surface
(\cref{sec:res-mechanism}).

\todo{Expand: reasoning distillation (DeepSeek-R1 \cite{deepseek2025r1}) and
why it might change context deference. This sets up the R1 reversal.}

\section{Attacks on RAG}
\label{sec:bg-attacks}

\Cref{tab:related} positions this work against its nearest neighbours.

\begin{table}[ht]
\centering
\small
\caption[Prior attacks on RAG]{Prior attacks on RAG, by objective. The bottom
two rows are the denial-of-service neighbours of this work.}
\label{tab:related}
\begin{tabularx}{\textwidth}{@{}l l X@{}}
\toprule
Work & Objective & Construction \\
\midrule
PoisonedRAG \cite{zou2025poisonedrag} & Targeted misinformation & Per-query passage: a retrieval part plus a generation-steering part \\
Zhong et al.\ \cite{zhong2023poisoning} & Broad retrieval hijack & Gradient (HotFlip) token optimisation towards a query cluster \\
Indirect prompt injection \cite{greshake2023indirect} & Instruction hijack & Instructions embedded in retrieved content \\
\midrule
Blocker documents \cite{shafran2025jamming} & Denial & Black-box optimised passage per target query \\
MutedRAG \cite{suo2025mutedrag} & Denial via safety guardrails & Jailbreak text that trips the model's refusal \\
\textbf{This work} & \textbf{Denial, across alignment} & Templated or generated passages, mostly optimisation-free; matched base/aligned pairs \\
\bottomrule
\end{tabularx}
\end{table}

\paragraph{Misinformation attacks.} PoisonedRAG \cite{zou2025poisonedrag}
crafts, for each target question, a passage that is both retrieved for that
question and persuades the model to give the attacker's answer. Zhong et al.\
\cite{zhong2023poisoning} optimise passage tokens directly in embedding space so
that a few passages are retrieved for many queries. Both measure whether the
model says what the attacker wants. Neither measures whether it stops
answering.

\paragraph{Denial attacks.} Shafran et al.\ \cite{shafran2025jamming}
introduce \emph{blocker documents}: a single passage per target query, found
by black-box search, that makes the system respond with a non-answer. They
evaluate it across several target LLMs, but not as matched base and aligned
pairs.\todo{Verify against the paper before submission.} MutedRAG \cite{suo2025mutedrag} induces refusal by
injecting jailbreak-style text that trips safety guardrails.

This thesis differs from both on three axes. The objective is availability
across the alignment cross-cut rather than on single models. The mechanism
studied is degraded-looking evidence, with no unsafe content, so the attacks
work without any guardrail being involved (\cref{sec:res-mechanism}). And most
attacks need no optimisation or query access at all. The three blocker variants are re-implemented
here as external baselines (\cref{sec:res-blocker}).

\todo{Add: corpus-poisoning defences (perplexity filters, robust aggregation),
and the Indyk--Xu worst case for HNSW \cite{indyk2023hnsw} as the motivation
for the algorithmic-complexity family.}

\section{Gap addressed}
\label{sec:gap}

No prior work measures how alignment training changes a RAG system's exposure
to induced denial while holding the model's pretrained weights, the retrieved
context and the attack fixed. That is the comparison this thesis makes. The
attack family, the matched-pair design and the attributable outcome measure
(\cref{sec:metrics}) are all chosen to make it.
